%% parts/p05/provability-truth.tex --- Part 5: Provability, the Second
%% Incompleteness Theorem, and Truth (Phase-3 pilot migration of v1.1
%% ch:second-incompleteness).  Body = statements only; proofs live in
%% annex/p05.tex; the v1.1 final audit table survives as the hand-authored
%% Table tab:p05-final-audit alongside the generated part register.
\chapter{Provability, the Second Incompleteness Theorem, and Truth}
\label{sec:p05-provability-truth}
\chaptermark{Provability, Incompleteness, and Truth}

\section{Derivability conditions}
\label{sec:p05-provability-truth-01}

Fix a computably enumerable theory \(T\supseteq\mathsf Q\) and use the
proof predicate and associated formula \(\mathrm{Prov}_T(x)\) fixed in
Definition~\ref{def:p05-proof-predicate}.  To reduce clutter, write
\(\mathrm{Prov}_T(\ulcorner\varphi\urcorner)\) for
\(\mathrm{Prov}_T(\overline{\ulcorner\varphi\urcorner})\).

\begin{hbdefinition}{def:p05-derivability-conditions}{Hilbert--Bernays--L\"ob derivability conditions}
The provability predicate satisfies the \emph{derivability conditions}
when, for all sentences \(\varphi\) and \(\psi\):
\begin{enumerate}
  \item if \(T\derives\varphi\), then
        \(T\derives\mathrm{Prov}_T(\ulcorner\varphi\urcorner)\);
  \item
  \[
    T\derives
    \mathrm{Prov}_T(\ulcorner\varphi\to\psi\urcorner)
    \to
    \bigl(\mathrm{Prov}_T(\ulcorner\varphi\urcorner)
    \to\mathrm{Prov}_T(\ulcorner\psi\urcorner)\bigr);
  \]
  \item
  \[
    T\derives
    \mathrm{Prov}_T(\ulcorner\varphi\urcorner)
    \to
    \mathrm{Prov}_T
    (\ulcorner\mathrm{Prov}_T(\ulcorner\varphi\urcorner)\urcorner).
  \]
\end{enumerate}
The first clause is a metatheoretic implication; the latter two are
theorem schemes internal to \(T\).  The three clauses are cited below as
D1, D2, and D3.
\end{hbdefinition}

For \(T=\PA\) and its fixed proof predicate, all three derivability
conditions hold: clause~1 is the displayed metatheoretic implication,
while \(\PA\) derives the schemes in clauses~2 and~3 by arithmetizing
proof concatenation and proof verification
\cite{HilbertBernays1939,Lob1955}.  For an arbitrary computably
enumerable \(T\supseteq\mathsf Q\), the results below state the three
conditions as hypotheses rather than inferring them from a vague strength
requirement.

\begin{hbdefinition}{def:p05-consistency-sentence}{Formal consistency sentence}
Fix \(\bottom_A\) as the false arithmetic sentence \(0=S0\).  Define
\[
  \mathsf{Con}(T)\defeq
  \neg\mathrm{Prov}_T(\overline{\ulcorner\bottom_A\urcorner}).
\]
This formula expresses consistency relative to the fixed numbering and
proof predicate.
\end{hbdefinition}

\begin{hbwarning}{def:p05-internal-consistency}{External and internal consistency}
The metalinguistic claim that \(T\) is consistent is not literally the
formula \(\mathsf{Con}(T)\).  Correctness of the chosen proof predicate
connects them in the standard structure \(\mathcal N_{\mathrm{ar}}\); the second
incompleteness theorem concerns whether \(T\) proves this particular
internal sentence.
\end{hbwarning}

\section{L\"ob's theorem}
\label{sec:p05-provability-truth-02}

\begin{hbtheorem}{thm:p05-lob}{proved=annex:p5.6}{Lob's theorem@L\"ob's theorem}
Let \(T\supseteq\mathsf Q\) be computably enumerable, and assume its
fixed provability predicate (Definition~\ref{def:p05-proof-predicate})
satisfies the three derivability conditions.  For every sentence
\(\varphi\), if
\[
  T\derives
  \mathrm{Prov}_T(\ulcorner\varphi\urcorner)\to\varphi,
\]
then \(T\derives\varphi\).
\end{hbtheorem}

\section{The second incompleteness theorem}
\label{sec:p05-provability-truth-03}

\begin{hbtheorem}{thm:p05-goedel-second}{proved=annex:p5.7}{Godel's second incompleteness theorem@G\"odel's second incompleteness theorem}
Let \(T\) be a consistent computably enumerable extension of
\(\mathsf Q\) whose provability predicate, fixed in
Definition~\ref{def:p05-proof-predicate}, satisfies the derivability
conditions.  Then
\[
  T\notderives\mathsf{Con}(T).
\]
\end{hbtheorem}

\begin{hbwarning}{def:p05-second-theorem-scope}{Scope of the second theorem}
The conclusion is not ``\(T\) cannot prove any sentence expressing any
notion of consistency.''  It concerns the arithmetized proof predicate
fixed in Definition~\ref{def:p05-proof-predicate} and the corresponding
sentence \(\mathsf{Con}(T)\).  Other deliberately weak predicates can
fail the derivability conditions.
\end{hbwarning}

\begin{hbexample}{def:p05-relative-consistency}{Relative consistency is not excluded}
A stronger theory \(T'\) may prove \(\mathsf{Con}(T)\); the second
theorem places no prohibition on such a relative result.  For example,
sufficiently strong standard set theories verify the existence of the
standard natural number structure and its satisfaction of the axioms of
\(\PA\), from which they formalize \(\mathsf{Con}(\PA)\).  This is a
claim about what the stronger theory proves, not a claim that its
consistency alone implies \(\mathsf{Con}(T)\).  The theorem says that a
consistent \(T\) cannot prove its own fixed consistency sentence under
the stated hypotheses.
\end{hbexample}

\section{Tarski's undefinability theorem}
\label{sec:p05-provability-truth-04}

\begin{hbdefinition}{def:p05-truth-definition}{Arithmetical truth definition}
An \(\Lang_A\)-formula \(\mathrm{Tr}(x)\) is an \emph{arithmetical truth
definition} if, for every \(\Lang_A\)-sentence \(\sigma\),
\[
  \mathcal N_{\mathrm{ar}}\entails
  \mathrm{Tr}(\overline{\ulcorner\sigma\urcorner})
  \leftrightarrow\sigma.
\]
The biconditional is evaluated in the standard structure and ranges
metalinguistically over all arithmetic sentences.
\end{hbdefinition}

\begin{hbtheorem}{thm:p05-tarski}{proved=annex:p5.8}{Tarski's undefinability theorem}
There is no arithmetical truth definition for the full set of true
\(\Lang_A\)-sentences.
\end{hbtheorem}

\begin{hbcorollary}{thm:p05-truth-undefinable}{proved=annex:p5.9}{Undefinability of the arithmetical truth set}
The set of G\"odel numbers
\[
  \{\ulcorner\sigma\urcorner:
  \sigma\text{ is an arithmetic sentence and }
  \mathcal N_{\mathrm{ar}}\entails\sigma\}
\]
is not definable in \(\mathcal N_{\mathrm{ar}}\) by any \(\Lang_A\)-formula.  In
particular, it is not computably enumerable.
\end{hbcorollary}

\begin{hbexample}{def:p05-partial-truth}{Partial truth definitions}
For each fixed level of the arithmetical hierarchy, arithmetic can define
a satisfaction predicate for formulas of that bounded complexity.
Tarski's theorem excludes one arithmetic formula correct for all
arithmetic sentences at once.  This distinction is essential in
reflection principles and proof theory \cite{Tarski1983}.
\end{hbexample}

\section{Audit of hypotheses}
\label{sec:p05-provability-truth-05}

Table~\ref{tab:p05-final-audit} is the hand-authored audit of the exact
hypotheses of this part's principal results, kept alongside the generated
register that closes the part; the third column names the invalid
strengthening most often substituted for each claim.

\begingroup\small
\begin{longtable}{@{}p{0.24\textwidth}p{0.32\textwidth}p{0.34\textwidth}@{}}
\caption{Exact hypotheses of the incompleteness and undefinability
  theorems, with common invalid strengthenings.}
\label{tab:p05-final-audit}\\
\toprule
Claim & Required hypotheses & Common invalid strengthening \\
\midrule
\endfirsthead
\toprule
Claim & Required hypotheses & Common invalid strengthening \\
\midrule
\endhead
First incompleteness, G\"odel sentence & \(T\) computably enumerable,
\(T\supseteq\mathsf Q\); consistency for unprovability; omega-consistency
or \(\Sigma^0_1\)-soundness for independence of \(G_T\), that is, for
underivability of both \(G_T\) and \(\neg G_T\) & Every consistent theory
has the displayed G\"odel sentence independent under no effectiveness or
arithmetic-strength assumption. \\
\addlinespace
First incompleteness, Rosser & \(T\) consistent, computably enumerable,
and \(T\supseteq\mathsf Q\) & Every consistent theory whatever, with no
computable-enumerability or \(\mathsf Q\)-extension hypothesis, has an
independent sentence. \\
\addlinespace
Second incompleteness & \(T\supseteq\mathsf Q\) consistent and computably
enumerable; the previously fixed provability predicate satisfies all
three derivability conditions &
No theory can prove a consistency statement about itself or any weaker
theory. \\
\addlinespace
L\"ob's theorem & Diagonal lemma and all three derivability conditions &
Any informal predicate called ``provable'' satisfies L\"ob's theorem. \\
\addlinespace
Tarski undefinability & One arithmetic formula purporting to define truth
for all arithmetic sentences in \(\mathcal N_{\mathrm{ar}}\) & No restricted
satisfaction or partial truth predicate can be defined. \\
\bottomrule
\end{longtable}
\endgroup
